\documentclass[a4paper]{article}
\input{../../../../../newpreamble.tex}
\begin{document}
\lhead{Math 220: Homework 5}
\chead{Lance Remigio}
\rhead{\thepage}
\begin{problem}
   Let \( X  \) and \( Y  \) be two topological spaces and \( f : X \to Y  \). 
   \begin{enumerate}
       \item[(i)] Prove that \( f  \) is continuous if and only if \( f^{-1}(C)  \) is closed for all closed sets \( C  \) in \( Y  \).
           \begin{proof}
               \( (\Longrightarrow) \) Suppose \( f  \) is continuous. Let \( C  \) be a closed set in \( Y  \). Our goal is to show that \( f^{-1}(C)  \) is closed in \( X  \). Since \( C  \) is closed in \( Y  \), it follows that its complement \( C^{c} \) is open in \( Y  \). Since \( f \) is continuous, it preimage \( f^{-1}(U) \) where \( U = C^{c} \) is open in \( X  \). This tells us that \( [f^{-1}(U)]^{c} \) is closed in \( X  \). Since \( [f^{-1}(U)]^{c} = f^{-1}(U^{c})  = f^{-1}(C) \), it follows that \( f^{-1}(C) \) is closed in \( X  \). 
\( (\Longleftarrow) \) Suppose that for all closed sets \( C  \) in \( Y  \), \( f^{-1}(C) \) is closed in \( X  \). Our goal is to show that \( f \) is continuous; that is, we want to show that for all open sets \( U  \) in \( Y  \), \( f^{-1}(U) \) is open in \( X  \). To this end, let \( U  \) be an open set in \( Y  \). This tells us that \( U^{c} = V  \) is a closed set in \( Y  \). Then \( f^{-1}(V) \) is closed in \( X  \) by assumption. That is, \( [f^{-1}(V)]^{c} \) is open in \( X  \). Since 
               \[  [f^{-1}(V)]^{c} = f^{-1}(V^{c}) = f^{-1}(U), \]
               it follows that \( f^{-1}(U) \) is open in \( X  \) which is our desired result.
           \end{proof}
       \item[(ii)] Prove that \( f  \) is continuous if and only if \( f(\overline{A})  \subseteq \overline{f(A)}   \) for all \( A \subseteq  X  \).
           \begin{proof}
               \( (\Longrightarrow) \)
               Note that \( \overline{f(A)} \) is a closed set in \( Y  \) and so its preimage under \( f  \), \( f^{-1}(\overline{f(A)}) \), is closed in \( X  \) from the previous part. Since \( f(A) \subseteq \overline{f(A)} \), we have that \( A \subseteq f^{-1}(\overline{f(A)}) \). Since \( f^{-1}(\overline{f(A)})\) is a closed set in \( X  \) containing \( A  \) and \( \overline{A} \) is the smallest closed set in \( X  \) containing \( A  \), we have that \( \overline{A} \subseteq f^{-1}(\overline{f(A)}) \), meaning that \( f(\overline{A}) \subseteq  \overline{f(A)} \).

                \( (\Longleftarrow) \) Let \( C  \) be a closed set in \( Y  \). Our goal is to show that \( f^{-1}(C) \) is closed in \( X  \); that is, we want to show that \( \overline{f^{-1}(C)} = f^{-1}(C) \). Notice that \( f^{-1}(C) \subseteq \overline{f^{-1}(C)} \) is immediate. So, it suffices to show that \( \overline{f^{-1}(C)} \subseteq f^{-1}(C) \). Indeed, by setting \( A = f^{-1}(C) \) which is a subset of \( X  \) and using our assumption, we get that 
                \begin{align*}
                    f(\overline{A}) \subseteq \overline{f(A)} &\implies \overline{A} \subseteq f^{-1}(\overline{f(A)}) \\
                                                              &\implies \overline{f^{-1}(C)} \subseteq f^{-1} \Big[ \overline{f (f^{-1}(C))}\Big] \\
                                                              &\implies \overline{f^{-1}(C)} \subseteq f^{-1}(\overline{C}) \\
                                                              &\implies \overline{f^{-1}(C)} \subseteq f^{-1}(C) \tag{\( \overline{C} = C \) since \( C  \) is closed in \( Y \)}
                \end{align*}
                Since \( \overline{f^{-1}(C)} \subseteq  f^{-1}(C) \) and \( f^{-1}(C) \subseteq \overline{f^{-1}(C)} \), we have \( \overline{f^{-1}(C)} = f^{-1}(C) \) and so we conclude that \( f^{-1}(C) \) is closed in \( X  \). Thus, \( f \) is continuous using part (i).
           \end{proof}
       \item[(iii)] Prove that if, \( f  \) is continuous and \( A  \) is dense in \( X  \), then \( f(A) \) is dense in \( f(X) \).
           \begin{proof}
           Our goal is to show that \( \overline{f(A)}^{f(X)} = f(X) \); that is, we want to show that \( f(X) \subseteq  \overline{f(A)}^{f(X)} \) and \( \overline{f(A)}^{f(X)} \subseteq f(X) \). Starting with the first inclusion, we observe that since \( f  \) is continuous and \( A  \) is dense in \( X  \) (i.e \( \overline{A} = X  \)), it follows from part (ii) that (using a proposition proven in class)
           \begin{align*}  
               f(X) = f(\overline{A}) \subseteq \overline{f(A)}^{Y} &\implies f(X) \subseteq \overline{f(A)}^{Y} \\ 
                                                                    &\implies f(X) \subseteq \overline{f(A)}^{Y} \cap f(X) \\
                                                                    &\implies f(X) \subseteq \overline{f(A)}^{f(X)}.
           \end{align*}
           Now, the other inclusion follows immediately because, via the same proposition stated before, we have  
           \[ \overline{f(A)}^{f(X)} = \overline{f(A)}^{Y} \cap f(X) \subseteq f(X).  \]
           Hence, it follows that \( \overline{f(A)}^{f(X)} = f(X) \) and we conclude that \( f(A) \) is dense in \( f(X) \).
           \end{proof}
   \end{enumerate}
\end{problem}

\begin{problem}[Continuity and Lower Limit Topology]
   \begin{enumerate}
       \item[(i)] Let \( {\mathcal{T}}_{L} \) be the lower limit topology and \( {\mathcal{T}}_{E} \) be the Euclidean topology on \( \R  \). Prove that \( f: (\R , {\mathcal{T}}_{L}) \to (\R , {\mathcal{T}}_{E}) \) is continuous if and only if for every \( x \in \R  \) and \( \epsilon > 0  \), there exists \( \delta > 0  \) such that \( | f(x) - f(y) |  < \epsilon  \) for all \( y \in [x,x+\delta) \).
           \begin{proof}
               \( (\Longrightarrow) \) Suppose \( f \) is continuous. Our goal is to show that for all \( x \in \R  \) and \( \epsilon > 0  \), there exists a \( \delta > 0  \) such that \( | f(x) - f(y) |  < \epsilon  \) for all \( y \in [x,x+\delta) \); that is, we need to show that for all \( x \in \R  \) and \( \epsilon > 0 \), there exists a \( \delta > 0 \) such that  
               \[  [x,x+\delta) \subseteq f^{-1}({B}_{\epsilon}(f(x))). \]
               Let \( x \in \R  \) and \( \epsilon > 0  \) be given. Now, consider the open ball \( {B}_{\epsilon}(f(x)) \) and notice that \( x \in f^{-1}({B}_{\epsilon}(f(x))) \). Since \( {B}_{\epsilon}(f(x)) \) is an open set in \( (\R, {\mathcal{T}}_{E}) \) and \( f  \) is continuous, we know that \( f^{-1}({B}_{\epsilon}(f(x))) \) is open in \( (\R , {\mathcal{T}}_{L}) \). Since \( x \in f^{-1}({B}_{\epsilon}(f(x)))  \), there exists a basis element \( U' \) containing \( x  \) in \( {\mathcal{T}}_{L} \) such that \( x \in U' \subseteq  f^{-1}({B}_{\epsilon}(f(x))) \). Since \( U' \) is a basis element, \( U' = [a,b) \) for some \( a,b \in \R \). Since \( x \in U' \), define \( \delta = \frac{ 1 }{ 2 }  (b - x) \) (this is strictly positive because \( x \in [a,b) \); that is, \( a \leq x < b \implies b - x > 0  \) ). Thus, we see that \( [x, x + \delta) \subseteq U' \subseteq f^{-1}({B}_{\epsilon}(f(x))) \) and so we conclude that  
               \[  [x,x+\delta) \subseteq f^{-1}({B}_{\epsilon}(f(x))). \]

               \( (\Longleftarrow) \) Conversely, we will show that \( f  \) is continuous. Suppose that for all \( x \in \R  \) and \( \epsilon > 0  \), there exists a \( \delta > 0  \) such that \( [x , x + \delta) \subseteq f^{-1}({B}_{\epsilon}(f(x))) \). To show that \( f  \) is continuous, we need to show that for all open sets \( U  \) in \( (\R, {\mathcal{T}}_{E}) \), \( f^{-1}(U) \) is open in \( (\R, {\mathcal{T}}_{L}) \). 

               To this end, let \( U  \) be an open set with respect to \( {\mathcal{T}}_{E} \). Let \( x \in f^{-1}(U) \) be arbitrary. Then \( f(x) \in U \), by definition. Since \( U  \) is open in \( {\mathcal{T}}_{E} \), there exists an \( \epsilon > 0  \) such that \( {B}_{\epsilon}(f(x)) \subseteq U  \) with respect to the standard metric on \( \R  \). By assumption, since \( x \in \R  \) and \( \epsilon > 0  \),
               \begin{center}
               there exists a \( \delta > 0  \) such that \( [x, x + \delta) \subseteq f^{-1}({B}_{\epsilon}(f(x))) \).  
               \end{center}
               Since \( [x,x+ \delta) \in {\mathcal{T}}_{L} \) (that is, \( [x, x + \delta) \) is open with respect to the Lower Limit Topology) contains \( x  \), it follows that 
               \[  x \in [x,x + \delta) \subseteq f^{-1}({B}_{\epsilon}(f(x)) ) \subseteq f^{-1}(U) \implies x \in [x , x + \delta) \subseteq f^{-1}(U). \]
               Thus, this implies that \( f^{-1}(U) \) is open in \( {\mathcal{T}}_{L} \) and so \( f \) is continuous.
           \end{proof}
       \item[(ii)] Is \( f : (\R, {\mathcal{T}}_{L}) \to (\R, \mathcal{T}_{E}) \), \( f(x) = \displaystyle 
       \begin{cases}
       1 & \text{if} \ x \geq 0 \\      
       0 & \text{if} \ x < 0
       \end{cases} \   \) continuous?
       \begin{solution}
           We claim that \( f \) is continuous. We will show that for all \( U \in \mathcal{T}_{E} \), the preimage under \( f^{-1}(U) \) is open with respect to the lower limit topology. Indeed, let \( U  \) be an open set in \( {\mathcal{T}}_{E} \). Since the image under \( f  \) only consists of either \( 1  \) or \(  0  \), we will consider four cases:
           \begin{enumerate}
               \item[(1)] Both \( 0  \) and \( 1  \) are not in \( U  \);
                \item[(2)] \( 0 \in U \) but \( 1 \notin U  \);
                \item[(3)] \( 0 \notin U  \) and \( 1 \in U  \);
                \item[(4)] \( 0 \in U  \) and \( 1 \in U \).
           \end{enumerate}
           Starting with (1), we notice that there is no \( x \in \R  \) such that \( f(x) \in U  \) and so \( f^{-1}(U) = \emptyset \) which is clearly in \( \mathcal{T}_{L} \).  
           
           Assuming (2), we notice that \( f(x) = 0  \) if and only if \( x < 0  \). Thus, the preimage of \( U  \) under \( f  \) is \( f^{-1}(U) = (-\infty , 0 ) \) which can be written as a union of basis elements of \( \mathcal{T}_{L} \) in the following way:
           \[  (-\infty , 0) = \bigcup_{ n = 1  }^{ \infty   } [-n, 0) \]
           which is clearly open since \( [-n,0) \) are open for all \( n \in \N  \) with respect to the Lower Limit Topology. Thus, \( f^{-1}(U) \) is open.
            
           Assuming (3), we notice \( f(x) = 1  \) if and only if \( x \geq 0  \). Thus, \( f^{-1}(U) = [0,\infty) \) which is in \( \mathcal{T}_L \). Thus, \( f^{-1}(U) \) is open and so \( f \) is continuous.

           Lastly, assume (4). We observe that since both \( 1  \) and \( 0  \) are in \( U  \), it follows that the preimage under \( f  \) of \( U  \) is just \( \R  \) itself which is contained in \( \mathcal{T}_L \). Thus, \( f^{-1}(U) = \R  \) is open in \( {\mathcal{T}}_{L} \).
           Thus, in all four cases, we get that \( f^{-1}(U) \) is open in \( \mathcal{T}_{L} \) and so we conclude that \( f \) is continuous.
       \end{solution}
   \item[(iii)] Is \( f: (\R, {\mathcal{T}}_{E}) \to (\R, {\mathcal{T}}_{L}) \), \( f(x) = 
   \begin{cases}
       1 &\text{if} \ x \geq 0 \\
       0 & \text{if} \ x < 0
   \end{cases}  \) \ continuous?
   \begin{solution}
       We claim that this function is not continuous. Take the open set \( [1,2) \) with respect to the Lower Limit Topology and observe that its pre-image under \( f  \) will be  
       \begin{align*}
           f^{-1}[1,2) &= \{ x \in \R : f(x) \in [1,2) \}  \\
                       &= \{ x \in \R : 1 \leq f(x) < 2  \}  \\
                       &= \{ x \in \R : x \geq 0  \} \\
                       &= [0,+\infty)
       \end{align*}
       which is not open with respect to the Euclidean Topology since \( 0  \) is not an interior point of \( [0,+\infty ) \).
   \end{solution}
   \end{enumerate} 
\end{problem}

\begin{problem}[Pull-back topology]
    Let \( X  \) be a set, \( (Y,\mathcal{T}) \) be a topological space, and \( f: X \to Y  \) be a function.
    \begin{enumerate}
        \item[(i)] Prove that \( f^{*} \mathcal{T} = \{ f^{-1}(U) : U \in \mathcal{T} \}   \) is a topology on \( X  \). (This is called the pull-back topology on \( X  \) induced by \( f \).)
            \begin{proof}
            Our goal is to show the following conditions: 
            \begin{enumerate}
                \item[(i)] \( \emptyset, X \in f^{*} \mathcal{T} \);
                \item[(ii)] For all \( {A}_{1}, {A}_{2} \in f^{*} \mathcal{T} \), \( {A}_{1} \cap {A}_{2} \in f^{*} \mathcal{T} \);
                \item[(iii)] If \( \{ {A}_{\alpha} \}_{\alpha \in \Omega}  \) is an arbitrary collection in \( f^{*} \mathcal{T} \), then \( \bigcup_{ \alpha \in \Omega }^{  }  {A}_{\alpha} \in f^{*} \mathcal{T} \).
            \end{enumerate}
            (i) Notice that \( f^{-1}(\emptyset) = \emptyset \) and \( f^{-1}(Y) = X  \) where \( \emptyset, Y \in \mathcal{T} \). Thus, it follows that \( \emptyset, X \in f^{*} \mathcal{T} \).

            (ii) Let \( {A}_{1}, {A}_{2} \in f^{*} \mathcal{T} \). Then by definition, we have \( {A}_{1} = f^{-1}({U}_{1}) \) and \( {A}_{2} = f^{-1}({U}_{2})  \) for some \( {U}_{1} \) and \( {U}_{2} \) in \( \mathcal{T} \). Then
            \[  {A}_{1} \cap {A}_{2} = f^{-1}({U}_{1}) \cap f^{-1}({U}_{2}) = f^{-1}({U}_{1} \cap {U}_{2}). \]
            Notice that since \( {U}_{1} \) and \( {U}_{2}  \) are in \( \mathcal{T} \) and \( \mathcal{T}  \) is a topology on \( Y \), it follows that \( {U}_{1} \cap {U}_{2} \in \mathcal{T} \). Thus, it follows that \( {A}_{1} \cap {A}_{2} \in f^{*} \mathcal{T} \).

            (iii) Let \( \{ {A}_{\alpha}  \}_{\alpha \in \Omega}\) be a collection of sets in \( f^{*} \mathcal{T} \). Then for all \( \alpha \in \Omega \), \( {A}_{\alpha} = f^{-1}({U}_{\alpha}) \) for some \( {U}_{\alpha} \in \mathcal{T} \). Then
            \[  \bigcup_{ \alpha \in \Omega  }^{  }  {A}_{\alpha} = \bigcup_{ \alpha \in \Omega  }^{  }  f^{-1} ({U}_{\alpha}) = f^{-1} \Big( \bigcup_{  \alpha \in \Omega }^{  }  {U}_{\alpha} \Big). \]
            Since \( \{ {U}_{\alpha} \}_{\alpha \in \Omega} \subseteq \mathcal{T} \) and \( \mathcal{T} \) is a topology, it follows that \( \bigcup_{ \alpha \in \Omega }^{  } {U}_{\alpha} \in \mathcal{T} \). This tells us that  
            \[  \bigcup_{ \alpha \in \Omega }^{  } {A}_{\alpha} \in f^{*} \mathcal{T}. \]
            Thus, from (i) through (iii), we conclude that \( f^{*} \mathcal{T} \) is, indeed, a topology on \( X  \).
            \end{proof}
        \item[(ii)] Prove that \( f: (X, f^{*} \mathcal{T}) \to (Y, \mathcal{T}) \) is continuous.
            \begin{proof}
            Let \( U  \) be an open set in \( (Y,\mathcal{T}) \) and let \( x \in f^{-1}(U) \) and so, \( f(x) \in U \). Since \( U \) is open in \( (Y, \mathcal{T})\), there exists an open set \( W  \) in \( (Y, \mathcal{T}) \) containing \( f(x)  \) such that \( f(x) \in W \subseteq U  \). But this tells us that \( x \in f^{-1}(W) \subseteq f^{-1}(U) \) where \( f^{-1}(W) \) is open in \( (X,f^{*}\mathcal{T}) \). Hence, it follows that \( f^{-1}(U) \) is an open set in \( (X, f^{*}\mathcal{T}) \). Thus, \( f  \) is continuous.
            \end{proof}
        \item[(iii)] Prove that, if \( S  \) is a topology on \( X  \) such that \( f: (X,S) \to (Y,\mathcal{T}) \) is continuous, then \( f^{*} \mathcal{T} \subseteq S \).
            \begin{proof}
            Let \( A \in f^{*} \mathcal{T}   \) be arbitrary. Then it follows by definition that there exists an open set \( U \) in \( (Y,\mathcal{T}) \) such that \( A = f^{-1}(U) \). Since \( f: (X,S) \to (Y,\mathcal{T}) \) is continuous and \( U \) is an open set in \( (Y, \mathcal{T}) \), \( f^{-1}(U) \) is open in \( (X,S) \), but this tells us that \( A  \) is open with respect to \( S  \) in \( X  \); that is, \( A \in S  \). Hence, \( f^{*} \mathcal{T} \subseteq S \).
            \end{proof}
    \end{enumerate}
\end{problem}

\begin{remark}
   In the next problem, whenever I mention an open set in \( \R^{n}  \) where \( n \in \N \), I am assuming with respect to the Euclidean Topology on \( \R^{n} \). 
\end{remark}

\begin{problem}[Open and Closed Sets Via Continuous Functions]
   \begin{enumerate}
       \item[(i)] Prove that \( \{ x \in \R : \cos x > \sin x  \}   \) is open in \( \R \).
           \begin{proof}
           Define \( f(x) = \cos  x - \sin x  \) where \( f: \R \to \R \). Since \( \sin x  \) and \( \cos x  \) are continuous on \( \R \) and that the difference of continuous functions is continuous, it follows that \( f(x) \) is continuous. Consider the open set \( (0,+\infty ) \) in \( \R  \). We see that 
           \begin{align*}
               f^{-1}(0, +\infty ) &= \{ x \in \R : f(x) \in (0, +\infty ) \}  \\
                                   &= \{ x \in \R : f(x) > 0  \}  \\
                                   &= \{ x \in \R : \cos x - \sin x > 0  \} \\
                                   &= \{ x \in \R : \cos x > \sin x  \}.
           \end{align*}
           Since \( f \) is continuous and \( (0,+\infty ) \) is open in \( \R \), we know that \( f^{-1}((0,+\infty )) \) is open in \( \R \). That is, \( x  \in \R : \cos x > \sin x \) is open in \( \R \).
           \end{proof}
        \item[(ii)] Show that \( \{ (x,y) \in \R^{2} : x > y , \ x^{2} + y^{2} < 1  \}  \) is open in \( \R^{2} \).
            \begin{proof}
                First, note that the set above can be written as an intersection of the following two sets
                \begin{enumerate}
                    \item[(a)] \( \{ (x,y) \in \R^{2} : x > y \}  \),
                    \item[(b)] \( \{ (x,y) \in \R^{2} : x^{2} + y^{2} < 1  \}  \).
                \end{enumerate}
                Thus, our goal is to show that the sets in (a) and (b) are open so that we can conclude that their intersection (because a finite intersection of open sets is open) is open.

                (a) Define \( f(x,y) = x - y  \) where \( f: \R^{2} \to \R \). Since \( x  \) and \( y  \) are polynomials in \( \R  \) which are continuous, it follows that their difference \( f(x,y) \) is also continuous. Consider the open set \( (0,+\infty ) \) in \( \R \). Then it follows that 
                \begin{align*}
                   f^{-1}(0,+\infty )  &= \{ (x,y) \in \R^{2} : f(x,y) \in (0,+\infty ) \}  \\
                                       &= \{ (x,y) \in \R^{2} : f(x,y) > 0 \}  \\
                                       &= \{ (x,y) \in \R^{2} : x - y > 0 \} \\
                                       &= \{ (x,y) \in \R^{2} : x > y  \}.
                \end{align*}
            Since \( f \) is continuous and \( (0,+\infty ) \) is open in \( \R \), it follows that \( f^{-1}(0,+\infty ) \) is open in \( \R^{2} \). That is, \( \{ (x,y) \in \R^{2} : x > y \}   \) is open in \( \R^{2} \). 

                (b) Define \( f(x,y) = x^{2} + y^{2} \) where \( f: \R^{2} \to \R \). Notice that \( f \) is continuous because \( x^{2} \) and \( y^{2} \) are polynomials in \( \R \) which are continuous and that the sum of these two continuous functions is continuous. Thus, \( f \) is continuous on \( \R^{2} \). Now, consider the open set \( (-\infty , 1 )  \) in \( \R \). Since \( f  \) is continuous, we know that its preimage \( f^{-1}(-\infty , 1 ) \) is open in \( \R^{2} \). Then
            \begin{align*}
                f^{-1}(-\infty , 1) &= \{ (x,y) \in \R^{2} : f(x,y) \in (-\infty , 1) \}  \\
                                    &= \{ (x,y) \in \R^{2} : f(x,y) < 1  \}  \\
                                    &= \{ (x,y) \in \R^{2} : x^{2} + y^{2} < 1  \}.
            \end{align*}
            Hence, \( \{ (x,y) \in \R^{2} : x^{2} + y^{2} < 1  \}  \) is open in \( \R^{2} \).
 
            Now, the intersection of the sets in (a) and (b) are open and so the original set \( \{ (x,y) \in \R^{2} : x > y , x^{2} + y^{2} < 1  \}  \) are open in \( \R^{2} \). 
            \end{proof}
        \item[(iii)] Prove that \( \{ (x,y,z) \in \R^{3} : 4 \leq  x^{2} + y^{2} + z^{2} \leq 6 \}  \) is closed in \( \R^{3} \).
            \begin{proof}
                Define \( f: \R^{3} \to \R  \) by \( f(x,y,z) = x^{2} + y^{2} + z^{2} \). Since \( x^{2}, y^{2},  \) and \( z^{2} \) are polynomials in \( \R  \) which are continuous, it follows from the fact that a sum of continuous functions that \( f  \) is continuous. Consider the closed set \( [4,6] \) in \( \R  \). Since \( f \) is continuous, its preimage \( f^{-1}([4,6]) \) is closed in \( \R^{3} \). That is, 
                \begin{align*}
                    f^{-1} ([4,6]) &= \{ (x,y,z) \in \R^{3} : f(x,y,z) \in [4,6] \}  \\
                                   &= \{ (x,y,z) \in \R^{3} : 4 \leq f(x,y,z) \leq 6 \} \\
                                   &= \{ (x,y,z) \in \R^{3} : 4 \leq x^{2} + y^{2} + z^{2} \leq 6 \}.
                \end{align*}
                Hence, it follows that \( \{ (x,y,z) \in \R^{3}: 4 \leq x^{2} + y^{2} + z^{2} \leq 6  \}   \) is closed in \( \R^{3} \).
            \end{proof}
   \end{enumerate} 
\end{problem}

\begin{problem}[Pasting Lemma For A Finite Cover]
   Assume the hypothesis of the pasting lemma. Show that if \( {A}_{i} \) is closed for all \( i \in I  \) and \( I  \) is finite, then \( f  \) is continuous. 
\end{problem}
\begin{proof}
    Let \( U  \) be closed in \( Y  \) and let \( \{ {A}_{i} \}_{i \in I} \) be a finite cover of \( X  \); that is,
            \[  \bigcup_{ i = 1  }^{ n }  {A}_{i} = X. \]
    Our goal is to show that \( f^{-1}(U) \) is closed in \( X  \). Notice that 
\begin{align*}
    f^{-1}(U) = f^{-1}(U) \cap X &= f^{-1}(U) \cap \Big( \bigcup_{ i=1  }^{ n }  {A}_{i} \Big)  \\
                                 &= \bigcup_{ i= 1  }^{ n  }  [f^{-1}(U) \cap {A}_{i}].
\end{align*}
Note that since \( f(x) = {f}_{i}(x)  \) for all \( x \in {A}_{i} \), we get that for all \( 1 \leq i \leq n  \)
\begin{align*}
    {f}_{i}^{-1}(U) &= \{ x \in {A}_{i} : {f}_{i}(x) \in U  \}  \\
                    &= \{ x \in {A}_{i} : f(x) \in U  \}  \\
                    &= \{ x \in {A}_{i} : x \in f^{-1}(U) \}  \\
                    &= {A}_{i} \cap f^{-1}(U).
\end{align*}
Then we get that 
\[  f^{-1}(U) = \bigcup_{ i=1  }^{ n }  f^{-1}_i (U). \]
Since each \( {f}_{i} \) is continuous on each \( {A}_{i} \) and \( U \) is a closed set in \( Y \),  it follows that \( {f}_{i}^{-1}(U) \) is closed in \( {A}_{i} \). Since each \( {A}_{i} \) is closed in \( X  \), it follows that \( {f}_{i}^{-1}(U) \) is closed in \( X  \) for all \( 1 \leq i \leq n  \). Since a finite union of closed sets is closed, it follows that 
\[  f^{-1}(U) = \bigcup_{ i= 1  }^{ n }  f^{-1}_i(U) \]
is a closed set in \( X  \) and so we conclude that \( f \) is continuous.
\end{proof}

\begin{problem}
   \begin{enumerate}
       \item[(i)] Let \( X  \) be a topological space and \( A  \) is a subspace of \( X  \). Consider the inclusion \( i : A \to X  \), \( i(x) = x  \). Then
           \begin{enumerate}
               \item[(a)] \( A  \) is open in \( X \) if and only if the inclusion \( i \) is open.
                   \begin{proof}
                       Let \( \mathcal{T} \) be a topology on \( X  \) and let \( \mathcal{T} |_{A} \) be the subspace topology endowed from \( X  \). 
                      
                       \( (\Longrightarrow) \) Suppose \( A  \) is open. Our goal is to show that the inclusion map \( i  \) is open; that is, \( i(U) \in \mathcal{T} \) for all \( U \in \mathcal{T} |_{A}  \). Let \( U \in \mathcal{T} |_{A} \) be arbitrary (that is, \( U  \) is open in \( A  \)). We want to show that \( i(U)  \) is open in \( X  \). Since \( U \) is open in \( A  \), there exists an open set \( V  \) in \( X  \) such that \( U = V \cap A  \). Since \( A  \) and \( V \) are open in \( X  \), it follows that \( U  \) must be open because a finite intersection of open sets is open. Thus, \( i(U) = U  \) implies that \( i(U) \) must be open in \( X  \). 

                       \( (\Longleftarrow) \) Suppose that \( i  \) is open. Then \( i \) is open if and only if for all \( U \in \mathcal{T} |_{A}  \), \( i(U) \in \mathcal{T} \). In particular, \( A \in \mathcal{T} |_{A} \), and so \( i(A) \in \mathcal{T} \). But note that \( i(A) = A  \) and so \( A \in \mathcal{T} \). Thus, it follows that \( A  \) is open in \( X  \).                    
                   \end{proof}
                \item[(b)] \( A  \) is closed in \( X  \) if and only if the inclusion \( i \) is closed.
                    \begin{proof}
                        \( (\Longrightarrow) \) Suppose \( A  \) is closed in \( X  \). Let \( U  \) be a closed set in \( A  \). Then there exists a closed set \( V  \) in \( X  \) such that \( U = V \cap A  \). However, \( A \) and \( V  \) are both closed in \( X  \) which implies that \( U = V \cap A   \) is closed because the finite intersection of closed sets is closed. Thus, \( i(U) =  U \) is closed in \( X  \) which is our desired result.                         

                        \( (\Longleftarrow) \) Suppose that \( i  \) is closed. That is, for all closed sets \( V \) in \( A  \), \( i(V) \) is closed in \( X  \). Since \( A  \) itself is closed in its subspace topology, it follows that \( i(A) \) is closed in \( X  \). But by definition \( i(A) = A  \) and so \( A  \) is closed in \( X  \).

                    \end{proof}
           \end{enumerate}
       \item[(ii)] Let \( n,k \in \N  \). Prove that the function
           \[  f: \R^{n} \to \R^{n+k} \ \ f({x}_{1}, {x}_{2}, \dots, {x}_{n}) = ({x}_{1}, {x}_{2}, \dots, {x}_{n}, 0 , 0 , \dots, 0) \]
           is not open.
           \begin{proof}
               Equip \( \R^{n} \) and \( \R^{n+k} \) with the Euclidean Topology. Since \( \R^{n}  \) and \( \R^{n+k} \) are metrizable topological spaces, we can equip both of them with the Euclidean norm \( \|\cdot\| \). Define the open ball \( U = {B}_{\epsilon}(0) \) centered at \( 0  \) with an arbitrary radius \( \epsilon > 0   \) in \( \R^{n} \) which is clearly open with respect to the Euclidean Topology on \( \R^{n} \). The image of this ball under \( f  \) is \[  f(U) = \{ ({x}_{1}, {x}_{2}, \dots, {x}_{n}, 0, 0, \dots , 0 ) \in \R^{n+k} : ({x}_{1}, {x}_{2}, \dots, {x}_{n}) \in {B}_{\epsilon}(0) \}. \] 
              We claim that \( 0 \in f(U) \) is not an interior point of \( f(U) \), implying that \( f(U)  \) is not open in \( \R^{n+k} \). Let \( r > 0  \), and consider the ball \( {B}_{r}(0) \subseteq \R^{n+k} \). Notice that this ball contains points that have at least one \( i \)-th component where \( i > n  \) that is non-zero. In particular, take the following point  
               \[  y = (0, \dots, 0, \dots, 0, r/2, 0, \dots, 0) \in {B}_{r}(0). \]
            We see that this point is contained in \( {B}_{r}(0) \) since  
            \[  \| y - 0  \| = \Big( \sum_{ i=1  }^{ n + k } | {y}_{i} - 0 |^{2}  \Big)^{1/2} = \Big| \frac{ r }{ 2 }  -  0  \Big| = \frac{ r }{ 2 }  < r   \]
            but not contained in \( f(U) \) because it contains some \( k  \)-th coordinate that is non-zero where \( k > n  \) (namely, from the above \( {y}_{k} = \frac{ r }{ 2 }  \)).      
               Hence, there does not exists an open ball centered at \(  0 \) that is contained in \( f(U) \). Therefore, \( f(U) \) is not open in \( \R^{n+k} \) and so \( f \) is not an open map. 
    \end{proof}
   \end{enumerate} 
\end{problem}


\end{document}

